% !TEX root = ../main.tex
% !TEX program = pdflatex
\subsection{Оценка эффективной скорости обработки массива}

Дополнительно оценена скорость последовательного цикла из метода A. Для каждого шага вычислялось отношение размера массива $S$ ко времени прохода $t$:
\[
  B_{\mathrm{eff}} = \frac{S}{t}.
\]
В расчёте $S=256$~МиБ для обоих шагов. При этом увеличение шага меняет число и расположение модифицируемых элементов. Поэтому величина характеризует эффективную скорость выбранного цикла по условному объёму массива, а не физический трафик DRAM.

По минимальным временам трёх запусков при шаге 64~байта получено 6,533--6,600~мс, что соответствует примерно 40,7--41,1~ГБ/с в десятичных единицах. При шаге 128~байт времена составили 8,115--8,131~мс, или около 33,0--33,1~ГБ/с. Эти значения относятся к циклу чтения с модификацией элементов; их нельзя интерпретировать как предельную пропускную способность памяти.
